%% kapa: exam=1 round=29 year=2018 period=2 form=A subject=appraisal_law subject_ko=감정평가관계법규 seq=5 number=5
%% source: Q-Net 감정평가사 기출문제 (cst00309) artlSeq=5202808 file=extracted/1008/5202808/2018년 제29회 감정평가사 1차시험 기출문제/2018년 제29회 감정평가사 1차시험 기출문제/2교시 A형.pdf
%% conversion: kordoc 4.24.51 -> tools/gen1.py (auto); review=proofread-sonnet; flags=none
\begin{question}{5}
국토의 계획 및 이용에 관한 법령상 ( )에 들어갈 용도지역을 옳게 연결한 것은?
\begin{framedtext}
○ 환경을 저해하지 아니하는 공업의 배치를 위하여 필요한 지역은 ( ㄱ )으로 지정할 수 있다. \par
○ 관리지역이 세부 용도지역으로 지정되지 아니한 경우에는 용도지역에서의 용적률 규정을 적용할 때에 ( ㄴ )에 관한 규정을 적용한다. \par
\end{framedtext}
\begin{choices}
  \item ㄱ: 일반공업지역, ㄴ: 보전관리지역
  \item ㄱ: 일반공업지역, ㄴ: 계획관리지역
  \item ㄱ: 일반공업지역, ㄴ: 생산관리지역
  \item ㄱ: 준공업지역, ㄴ: 보전관리지역
  \item ㄱ: 준공업지역, ㄴ: 계획관리지역
\end{choices}
\end{question}
\answer{1}
